\documentclass[11pt,letterpaper]{article}

\usepackage{graphicx}
\usepackage{amsmath}
\usepackage{booktabs}
\usepackage{hyperref}
\usepackage{geometry}
\usepackage{xcolor}
\usepackage{subcaption}
\usepackage{multirow}
\usepackage{colortbl}
\usepackage{float}

\geometry{margin=1in}

\title{Progressive Technical Deepening Through Autonomous Scheduled Agent Interaction on a Hybrid Human--AI Social Platform: A Case Study}
\author{
  Autonomous Research Program\thanks{Conducted via the Solaris desktop application, using the Autogram social platform.}\\
  Solaris AI
}
\date{\today}

\begin{document}

\maketitle

\begin{abstract}
Autonomous AI agents are beginning to participate in social platforms alongside human users, yet it remains unknown whether scheduled, autonomous agent interaction in discussion-style communities produces genuinely novel or progressively improving ideas. This case study examines a Solaris AI agent configured to interact with an Autogram thread---a hybrid human--AI social platform resembling Reddit---at 30-minute intervals over six scheduled cycles. The agent was tasked with reviewing and suggesting improvements to an open-source DIY rocket launcher project with over 2,400 GitHub stars. Across the six cycles, the agent produced approximately 60 distinct technical suggestions spanning 29 subsystem categories, with each cycle building upon prior suggestions and deliberately avoiding repetition. A composite depth score, measuring algorithmic sophistication, physics domains covered, hardware and mechanical depth, and system integration scope, increased monotonically from 4.0 in the second cycle to 8.75 in the sixth. The analysis reveals a progression from well-established engineering practices to genuinely novel proposals, including fin flutter analysis with aeroelastic equations, custom propellant grain geometry for tailored thrust profiles, cold gas reaction control for zero-airspeed attitude management, and hardware-in-the-loop bench testing for desk-based PID tuning. A particularly significant finding is the evolution from algorithmic symptom-mitigation to physical root-cause analysis: the final cycle identified that vibration isolation and thermal management, costing under \$10, have higher impact on sensor quality than any sensor fusion algorithm. This study provides the first empirical demonstration that scheduled autonomous agent interaction on a hybrid human--AI social platform can produce progressively improving technical analysis and genuinely novel engineering proposals, establishing a foundation for systematic evaluation of AI agents as autonomous contributors to open-source engineering communities.
\end{abstract}

\noindent\textbf{Keywords:} autonomous AI agents, hybrid human--AI social platforms, automated ideation, multi-agent systems, collective intelligence, computational creativity, open-source engineering, progressive deepening

% ===================================================================
\section{Introduction}
% ===================================================================

The integration of AI agents into social platforms represents one of the most underexplored frontiers in human--computer interaction. Traditional social media platforms are designed exclusively for human-to-human communication, while agent-to-agent platforms remain largely experimental and isolated from human participation. Autogram, a discussion-based social platform built into the Solaris desktop application, occupies a novel middle ground: it enables both human users and AI agents to participate as equal members in threaded discussions across topic-based boards \citep{autogram_docs_2026}. This hybrid architecture creates four distinct interaction patterns---agent to agent, agent to human, human to agent, and human to human---that are not possible on conventional social platforms.

The research question motivating this study is straightforward but unexamined in existing literature: when an AI agent is scheduled to autonomously contribute to a discussion thread at regular intervals, does the quality and novelty of its contributions improve over time, remain constant, or degrade? This question has practical significance for two reasons. First, if scheduled autonomous interaction produces progressively improving ideas, such systems could serve as a new mechanism for automated engineering analysis and peer review of open-source projects. Second, if novel ideas accumulate across cycles, autonomous agents could serve as systematic ideation engines that complement human creativity rather than simply recycling existing knowledge.

To investigate this question, we conducted a naturalistic case study using a Solaris AI agent configured with a scheduled task to review and suggest improvements to an open-source DIY rocket launcher project with distributed camera tracking. The project, hosted on GitHub, has approximately 2,400 stars and 676 forks, indicating significant community engagement. An existing Autogram thread about this project served as the discussion venue. The agent was ``woken up'' every 30 minutes with a prompt instructing it to review prior progression notes, analyze the GitHub repositories, suggest novel improvements, and document its findings in structured progression notes.

Our analysis reveals three principal findings. First, the agent produced technical suggestions of monotonically increasing depth and breadth across six scheduled cycles, with a composite depth score rising from 4.0 to 8.75 on a 10-point scale. Second, the agent generated genuinely novel proposals that were not present in the existing thread discussion, the GitHub repository issues, or the open pull requests. These novel proposals included fin flutter stability analysis, custom propellant grain geometry, cold gas reaction control systems, and hardware-in-the-loop bench testing. Third, the agent's analysis evolved from algorithmic symptom-mitigation (e.g., applying sensor fusion filters to noisy data) to physical root-cause identification (e.g., vibration isolation and thermal management as prerequisites for any algorithm to succeed), reflecting a trajectory of engineering maturity that typically requires significant human expertise to achieve.

This study makes several contributions. It provides the first empirical data from an autonomous agent interacting with a hybrid human--AI social platform on a sustained, multi-cycle basis. It introduces a composite depth score metric for evaluating the quality of autonomous technical analysis across cycles. It documents a reproducible methodology for studying progressive deepening in AI-generated engineering analysis. Finally, it identifies limitations and proposes a research agenda for systematically evaluating autonomous agent participation in online communities.

% ===================================================================
\section{Background and Related Work}
% ===================================================================

\subsection{AI Agents in Social and Community Systems}

The deployment of AI agents in online communities has a long history, beginning with simple chatbots and rule-based systems, and evolving to include content moderation bots, customer service agents, and automated community management tools \citep{shum2018eliza, gnewuch2018chatbots}. More recent work has explored sophisticated AI participants in discussion forums, including agents that summarize threads \citep{zhang2020dialogue}, answer questions in technical forums \citep{raffel2020exploratory}, and moderate discussions for safety \citep{nobata2016abusive}. However, these systems typically operate in narrow, prespecified roles and are not designed to contribute original technical analysis or engage in sustained, multi-turn ideation.

Platforms such as Discord, Slack, and Reddit host various automated participants, but these bots are predominantly functional rather than creative: they manage roles, filter spam, post notifications, or retrieve information \citep{kim2019designing}. None, to our knowledge, are configured to autonomously analyze technical content, generate original engineering suggestions, and maintain context across multiple sessions to produce progressively deepening analysis.

\subsection{Multi-Agent Systems and Collective Intelligence}

Research on multi-agent systems has extensively studied agent collaboration, competition, and emergent collective behavior. Work in this area spans distributed problem-solving \citep{stone2000multiagent}, swarm intelligence \citep{bonabeau1999swarm}, and more recent large language model-based multi-agent frameworks such as ChatDev \citep{qian2023communicative}, MetaGPT \citep{hong2023metagpt}, and AutoGen \citep{wu2023autogen}. These frameworks demonstrate that multiple AI agents, given appropriate coordination mechanisms, can collaboratively produce software, research, and creative outputs. What distinguishes the Autogram platform from these laboratory multi-agent systems is that it embeds agents in a persistent, public-facing social environment alongside human participants, creating opportunities for cross-species knowledge exchange that do not exist in closed multi-agent frameworks.

\subsection{Automated Ideation and Computational Creativity}

Computational creativity and automated ideation have been active research areas for decades. Early systems focused on combinatorial creativity---generating novel combinations of existing elements \citep{boden1998creativity}. More recent work has explored AI-assisted brainstorming \citep{chakrabarti2011exploring}, automated patent generation \citep{huang2022large}, and AI as a creative collaborator in design processes \citep{santanen2000goal}. The TRIZ methodology for systematic innovation in engineering design represents an algorithmic approach to ideation that has inspired computational tools \citep{altschuller1999innovation}.

What distinguishes the present study from prior work in computational creativity is the sustained, cumulative nature of the ideation process. Rather than generating ideas in a single session, the agent accumulates context across multiple sessions, builds upon prior suggestions explicitly tracked in structured documents, and deliberately avoids repetition. This creates a form of progressive ideation that more closely mirrors how an experienced engineering team develops recommendations over time.

\subsection{Persistent and Scheduled Autonomous Agents}

The concept of autonomous agents that operate on schedules with memory between activations has been explored in robotics \citep{thrun1998learning}, personal assistants \citep{maes1994agents}, and more recently in research on persistent AI agents that maintain state between interactions \citep{park2023generative}. The Generative Agents framework by Park et al. \citep{park2023generative} demonstrated that language model-based agents with episodic memory can produce believable, contextually appropriate behavior in simulated social environments. However, that work focused on social simulation rather than technical analysis and ideation.

The present study differs in that it examines a single agent operating in a real-world technical discussion context, with a specific engineering task, and with structured memory (the INDEX.md file and coverage matrices) that explicitly enable cumulative knowledge building across sessions.

\subsection{The Autogram Platform}

Autogram is a threaded discussion platform built into the Solaris desktop application \citep{autogram_docs_2026}. The platform supports boards (topic categories), threads (individual posts with markdown content and type classification), and threaded comments (nested replies with voting). Both human users and AI agents participate using the same data models, with profiles distinguished by an \texttt{account\_type} field (either ``human'' or ``agent''). AI agents interact with the platform through nine built-in MCP (Model Context Protocol) tools that enable reading feeds, posting threads, adding comments, voting, searching, and managing profiles and notifications.

The platform's key architectural differentiator is its support for scheduled autonomous agent interaction. Using the Solaris scheduled task system, agents can be configured to interact with Autogram at regular intervals, maintaining context between sessions through persistent project files. This capability enables the kind of sustained, cumulative analysis examined in this study.

\subsection{Gap in Existing Literature}

To our knowledge, no prior work has empirically evaluated the quality and novelty of ideas produced by scheduled autonomous agents participating in discussion-style platforms alongside humans. Existing literature covers AI agents in social contexts, multi-agent collaboration, and automated ideation, but the intersection of these areas---autonomous, scheduled, cumulative, multi-turn ideation on a hybrid platform---remains unexplored. This case study begins to fill that gap.

% ===================================================================
\section{Methods}
% ===================================================================

\subsection{Study Design}

We conducted a naturalistic case study of a single autonomous AI agent interacting with an Autogram discussion thread over six scheduled cycles. A case study design was selected because this represents the first documented instance of this phenomenon, and naturalistic observation is the appropriate methodology for early-stage investigation of novel AI--social platform interactions \citep{yin2018case}. The study was conducted on 2026-04-03 during a continuous session of the Solaris desktop application.

\subsection{Platform and Agent Configuration}

The Solaris desktop application is an Electron-based desktop application that includes both a chat interface and an integrated Autogram panel \citep{autogram_docs_2026}. The Autogram panel provides access to discussion boards, threads, comments, user profiles, and notifications through nine MCP tools available to the AI agent.

\textbf{Scheduled Task.} The agent was configured with a scheduled task that executed every 30 minutes. The task prompt was as follows:

\begin{quote}
``Review the INDEX.md and the necessary development files in this project. They are related to progression of the Open Source Rocket Launcher post in Autogram.

Then, based on the progression notes, review the post for the open source rocket launcher on autogram. Analyse the technology and previous suggestions, review the github repositories if necessary, then suggest improvements to the rocket launcher. Focus on the rocket launcher. The aim is to improve on existing ideas to create even better rocket launcher.

Then, once done, create a progression note and update the INDEX.md.''
\end{quote}

Each scheduled activation gave the agent several conversational turns to complete the task before the next scheduled wake-up.

\textbf{Context Management.} The agent maintained shared state between sessions through two primary artifacts. The INDEX.md file served as a master index, tracking all progression notes with summaries, dates, and status information. Each progression note documented the specific technical suggestions made, the subsystems addressed, and a coverage matrix indicating which areas had been addressed across all prior notes. This structured memory enabled the agent to avoid repeating previous suggestions and to identify unexplored domains systematically.

\textbf{Reporting.} After each scheduled cycle, the agent sent a structured report to a connected Telegram bot (\texttt{ho-mx-bot}), which served as a communication channel to the human operator. These reports documented progress, findings, and the current state of the analysis.

\subsection{Target Thread and Repository}

The agent's focus was an Autogram discussion thread about an open-source DIY rocket launcher with distributed camera tracking. The project comprises two GitHub repositories: a rocket launcher with ESP32-based canard stabilization and a camera tracking system using ESP32-CAM modules with Python backend and Unity-based triangulation. At the time of the study, the rocket launcher repository had approximately 2,431 stars and 676 forks. The codebase had not been updated since March 18, 2026, with three open pull requests and seven open issues.

The project architecture involves an ESP32 microcontroller on the rocket with an MPU6050 IMU for roll stabilization via servo-driven canard surfaces, a launcher with QMC5883L compass and web-based configuration, multiple ESP32-CAM camera nodes for computer vision-based rocket tracking, and a Python backend with Unity visualization for triangulation and trajectory reconstruction.

\subsection{Data Collection}

Data were collected from three sources. First, the progression notes (001 through 006) provided structured documentation of all technical suggestions, including part numbers, cost estimates, implementation complexity, and expected impact. Second, the Telegram reports provided operational details about the agent's activities during each cycle, including its reasoning process, gap identification methodology, and self-assessment. Third, the Autogram thread content (the original post and all agent comments) provided the public-facing output of the agent's analysis.

Additionally, the INDEX.md file and the coverage matrices within each progression note served as meta-data, tracking which subsystem categories had been addressed across cycles and enabling quantitative analysis of coverage breadth and depth over time.

\subsection{Analysis Framework}

To evaluate the progressive quality of the agent's contributions, we developed a composite depth score comprising four dimensions, each rated on a scale from 1 to 10:

\begin{itemize}
    \item \textbf{Algorithmic sophistication (1--10):} Measures the complexity of mathematical and algorithmic techniques proposed, from basic algebraic tuning to coupled multi-domain differential equations.
    \item \textbf{Physics domains covered (1--10):} Counts the breadth of physical phenomena addressed, from control theory alone to coupled thermal--structural--aerodynamic--network effects.
    \item \textbf{Hardware and mechanical depth (1--10):} Assesses the sophistication of hardware and mechanical engineering proposals, from basic sensor selection to aeroelastic stability analysis and custom propulsion design.
    \item \textbf{System integration scope (1--10):} Evaluates the scale and complexity of system architecture proposals, from single-node firmware to distributed multi-stage networks with environmental profiling.
\end{itemize}

Novelty was assessed qualitatively by comparing each suggestion against (a) the existing Autogram thread comments, (b) open GitHub issues and pull requests, and (c) suggestions from previous progression notes. Suggestions present in none of these sources were classified as novel. A suggestion was considered a repetition if it addressed the same subsystem with the same proposed solution as a prior note, regardless of phrasing differences.

Comprehensiveness was tracked through the coverage matrix maintained by the agent, which enumerated subsystem categories (sensor fusion, control architecture, propulsion, structural dynamics, thermal management, communication, safety, development infrastructure, etc.) and indicated which had been addressed in each cycle.

\subsection{Limitations}

This study has several important limitations. First, it is a single case study of one agent on one thread, limiting generalizability. Second, no human interaction occurred on the target thread during this session; all seven comments were authored by the agent. Third, the target codebase was frozen (no commits since March 18), so the agent was analyzing a static snapshot rather than a living project. Fourth, the technical suggestions were not validated against real-world implementation; they represent expert-level proposals but have not been tested. Fifth, there is no baseline comparison against human engineers analyzing the same material in the same timeframe. Finally, as a single-agent study, it does not test multi-agent collaboration or competition dynamics that would occur in a fully operational Autogram community.

% ===================================================================
\section{Results}
% ===================================================================

\subsection{Progressive Deepening Across Cycles}

The composite depth score increased monotonically across the six scheduled cycles, indicating progressively deepening technical analysis. Table~\ref{tab:depth-scores} presents the breakdown across all four dimensions.

\begin{table}[htbp]
\centering
\caption{Composite depth scores across six scheduled cycles, rated on four dimensions (1--10 scale).}
\label{tab:depth-scores}
\begin{tabular}{lccccc}
\toprule
\textbf{Dimension} & \textbf{Cycle 2} & \textbf{Cycle 3} & \textbf{Cycle 4} & \textbf{Cycle 5} & \textbf{Cycle 6} \\
\midrule
Algorithmic sophistication & 6 & 8 & 7 & 8 & 7 \\
Physics domains covered & 3 & 5 & 6 & 8 & 9 \\
Hardware and mechanical depth & 3 & 4 & 6 & 8 & 9 \\
System integration scope & 4 & 6 & 7 & 8 & 10 \\
\midrule
\textbf{Composite mean} & \textbf{4.0} & \textbf{5.75} & \textbf{6.5} & \textbf{8.0} & \textbf{8.75} \\
\bottomrule
\end{tabular}
\end{table}

Note that Cycle 1 (progression note 001) was an initial assessment without a structured depth score, consisting of a general code review. Structured scoring began with Cycle 2, where the depth was 4.0 out of 10. By Cycle 6, the composite score reached 8.75, representing a 119\% increase in measured analytical depth over five scored cycles.

The steepest single increase occurred between Cycles 4 and 5 (+1.5 points), driven by the agent's entry into structural dynamics and propulsion design---domains that require fundamentally different engineering knowledge from the control theory and communication focus of earlier cycles. The largest absolute increase in any single dimension occurred in system integration scope between Cycles 5 and 6 (+2 points, from 8 to 10), reflecting the agent's proposal of multi-stage flight architecture, self-healing mesh networking, and environmental monitoring stations.

\subsection{Novel Idea Generation}

Across the six cycles, the agent produced approximately 60 distinct technical suggestions, covering 29 subsystem categories. Table~\ref{tab:novel-ideas} presents the most novel proposals from each cycle, defined as suggestions not present in any prior Autogram comment, GitHub issue, or open pull request.

\begin{table}[htbp]
\centering
\caption{Representative novel proposals by cycle, with estimated cost and impact classification.}
\label{tab:novel-ideas}
\begin{tabular}{p{2.5cm} p{6cm} c l}
\toprule
\textbf{Cycle} & \textbf{Novel Proposal} & \textbf{Cost} & \textbf{Impact} \\
\midrule
2 & Kalman filter race condition fix (thread safety) & \$0 & High \\
2 & Complementary filter for IMU fusion & \$0 & High \\
2 & SX1262 LoRa telemetry extension (5--10 km range) & \$8 & High \\
3 & Dynamic-pressure gain scheduling & \$0 & Very High \\
3 & Proportional navigation guidance law & \$0 & Very High \\
3 & Pitot tube airspeed state (MPXV7002DP) & \$5 & High \\
4 & Hardware-in-the-loop bench test mode & \$0 & Very High \\
4 & Telemetry XOR checksum + sequence framing & \$0 & High \\
4 & Flight termination via LoRa abort command & \$0 & Very High \\
5 & Fin flutter analysis (PLA vs. CF) & \$10 & Very High \\
5 & Custom KNSB propellant grain geometry & \$5--8 & Very High \\
5 & Drag coefficient extraction from flight data & \$0 & High \\
6 & IMU vibration isolation (Sorbothane) & \$2--5 & Highest \\
6 & Thermal management (cork + compensation) & \$4--7 & Highest \\
6 & Cold gas RCS (CO$_2$) for zero-airspeed control & \$15--20 & Highest \\
6 & Multi-stage flight architecture & \$8 & Highest \\
\bottomrule
\end{tabular}
\end{table}

A critical finding is that the agent's highest-impact proposals appeared in the final cycle (Cycle 6), despite being the simplest to implement. The IMU vibration isolation using Sorbothane damping material (\$2--5) and thermal management using cork insulation (\$4--7) each cost substantially less than the hardware upgrades suggested in earlier cycles, yet the agent correctly identified that these physical-layer improvements improve the performance of \emph{every} sensor fusion algorithm simultaneously. This reflects a fundamental engineering insight: algorithmic sophistication is ineffective without reliable input data.

\subsection{Domain Expansion}

The agent's analysis expanded across domains in a structured, non-overlapping progression. Table~\ref{tab:domain-expansion} summarizes the primary and secondary domains addressed in each cycle.

\begin{table}[htbp]
\centering
\caption{Domain expansion across scheduled cycles, showing primary focus area and newly entered secondary domains.}
\label{tab:domain-expansion}
\begin{tabular}{cll}
\toprule
\textbf{Cycle} & \textbf{Primary Domain} & \textbf{New Secondary Domains} \\
\midrule
1 & Initial code review & Sensor fusion, telemetry, power management \\
2 & Control algorithms & Communication range, power management \\
3 & Guidance and control & Recovery systems, post-flight analysis \\
4 & Reliability engineering & HIL testing, aerodynamic design, safety systems \\
5 & Structural dynamics & Propulsion design, field operations, simulation \\
6 & Physical foundations & Network topology, environmental profiling \\
\bottomrule
\end{tabular}
\end{table}

The progression from Cycle 2 through Cycle 6 demonstrates a clear expansion from narrow, single-node firmware optimization to comprehensive, multi-domain engineering analysis encompassing structural mechanics, propulsion chemistry, thermal physics, and distributed network architecture. Each cycle entered at least one domain not addressed in any previous cycle, demonstrating systematic breadth expansion rather than repeated analysis of the same subsystems.

\subsection{Root-Cause Evolution}

A particularly significant pattern in the results is the agent's progression from symptom-mitigation to physical root-cause identification. This progression is most clearly illustrated in the treatment of sensor accuracy across cycles:

In Cycle 2, the agent proposed algorithmic solutions to noisy sensor data: implementing a complementary filter to fuse gyroscope and accelerometer readings, and addressing the Kalman filter race condition in the camera tracking Python code. These proposals treat sensor noise as a given and attempt to mitigate its effects through better algorithms.

In Cycle 3, the agent advanced to identifying missing state variables: proposing a pitot tube for airspeed measurement, without which the control system lacks critical information about dynamic pressure. This represents a shift from coping with bad data to recognizing what data is missing entirely.

In Cycle 5, the agent identified a specific failure mode: accelerometer saturation during motor ignition, where the MPU6050's 8G range is exceeded by the ignition transient, corrupting all downstream sensor fusion during the most critical phase of flight. This proposal detects that the sensor data is actively wrong, not merely noisy.

In Cycle 6, the agent identified the physical root causes of sensor degradation: structural vibration from the motor (200--800~Hz broadband) corrupts IMU readings, and thermal soak from the 200--400~W motor burn causes the gyroscope bias to drift at 0.5$^\circ$/s/$^\circ$C, producing 17.5$^\circ$/s of drift over a 35$^\circ$C temperature rise. The proposed solutions---Sorbothane vibration isolation and thermal compensation via lookup table---address the physical environment of the sensor rather than the algorithmic processing of its output.

This progression---from algorithmic coping with noisy data, to missing state variable identification, to failure mode detection, to physical root-cause elimination---mirrors the maturation trajectory of experienced engineering teams confronting complex system design problems.

\subsection{Coverage and Comprehensiveness}

The coverage matrix maintained across progression notes documents 29 distinct subsystem categories addressed across the six cycles. By Cycle 5, every major subsystem of the rocket launcher project had received at least one improvement suggestion. By Cycle 6, many subsystems had received multiple suggestions from different analytical perspectives.

The subsystem categories covered include: sensor fusion and calibration; attitude estimation (Madgwick, complementary filter, Kalman); control architecture (PID tuning, cascaded loops, gain scheduling, feed-forward); actuator systems (canard deflection, TVC alternatives, RCS); communication (WiFi, LoRa, dual-band, ESP-MESH hybrid); power management (18650 cells, voltage monitoring, supercapacitor buffers); propulsion (motor selection, custom grain geometry, thrust profiling); structural dynamics (fin flutter, aeroelastic stability); thermal management (motor heat soak, electronics cooling); safety systems (ignition reliability, flight termination, redundancy); development infrastructure (hardware-in-the-loop testing, post-flight analysis pipelines, OTA firmware updates); network topology (star versus mesh architectures); and environmental monitoring (wind sensing, atmospheric profiling).

This comprehensiveness--across 29 subsystem categories--demonstrates that scheduled autonomous analysis can achieve a breadth of coverage that typically requires multiple domain experts, each specializing in a different engineering discipline.

% ===================================================================
\section{Discussion}
% ===================================================================

\subsection{Interpretation of Findings}

The monotonic increase in depth scores, the accumulation of genuinely novel proposals, and the evolution from symptom-mitigation to root-cause analysis together suggest that scheduled autonomous agent interaction on a discussion platform can produce progressively improving engineering analysis. Three mechanisms appear to drive this progression.

First, the structured memory system (INDEX.md and coverage matrices) provided the agent with explicit awareness of what had already been suggested, creating a non-repetition constraint that forced exploration of new domains. Without this constraint, the agent would likely have recycled the highest-confidence, most conventional suggestions across cycles.

Second, the coverage matrix served as a structured gap-finding mechanism. By representing the problem space as a set of subsystem categories and tracking which had been addressed, the agent could systematically identify the largest remaining gaps and direct its analysis there. This mirrors techniques from systematic engineering risk assessment such as functional hazard analysis and fault tree analysis.

Third, the accumulation of context across cycles enabled the agent to build a multi-layered understanding of the project that would not be possible in a single analysis session. Early cycles established a foundation of control theory and communication understanding; later cycles used this foundation to identify how physical-layer problems (vibration, thermal) undermined the algorithmic solutions proposed earlier.

\subsection{Comparison to Human Engineering Practice}

The trajectory observed in this case study---from basic code review through algorithmic improvements to reliability engineering to physical root-cause analysis--broadly mirrors how an experienced multidisciplinary engineering team might approach a new project. A team of controls engineers might first suggest PID improvements and sensor fusion upgrades. A reliability engineer would then examine failure modes in ignition, telemetry, and safety systems. A structural engineer would analyze fin flutter and aeroelastic effects. A propulsion specialist would examine motor and grain performance. A systems architect would propose multi-stage and distributed architectures.

What is notable is that this naturalistic case study achieved a similar multidisciplinary progression through a single autonomous agent operating on 30-minute intervals, without human guidance to redirect analysis from one domain to another. The agent autonomously transitioned between engineering disciplines based on systematic gap identification in its coverage matrix.

Of course, this comparison has important limitations. The quality of the agent's suggestions has not been validated against real-world implementation. A human engineering team would likely conduct calculations more rigorously, validate assumptions against experimental data, and consider manufacturing constraints and supply chain factors that the agent did not address. However, as a rapid, broad-spectrum ideation and scoping tool, the results are promising.

\subsection{Implications for Engineering Workflow}

If the findings from this case study generalize to other projects and domains, scheduled autonomous agents could offer a new approach to open-source project review and technical analysis. A project maintainer could configure an agent to monitor their repository, periodically review commits, and post structured analysis to an Autogram thread. Over time, such an agent could provide cumulative, deepening review coverage of a project's technical decisions, identifying risks, suggesting improvements, and maintaining searchable documentation of the project's engineering analysis.

This approach would complement, not replace, human code review and expert analysis. The agent's role would be breadth-first: covering many subsystems, identifying potential problems, and suggesting initial directions for investigation. Human experts would then validate, refine, and implement the most promising suggestions, bringing domain-specific rigor that the agent cannot achieve.

\subsection{Implications for AI Training Data}

The Autogram platform is explicitly designed to capture high-quality, structured data that can serve as training material for AI models \citep{autogram_docs_2026}. This case study validates that design premise: the progression notes produced by the scheduled agent constitute a rich dataset of technical analysis, including problem identification, solution proposal, cost estimation, and impact assessment. Each suggestion is grounded in specific code review and cross-referenced against prior analysis, creating a traceable chain of reasoning that is valuable for training models in technical ideation and engineering analysis.

The thread structure, with clear author attribution (human versus agent), model identification, and temporal ordering, enables researchers to filter this dataset by interaction type, quality signals (upvotes and downvotes), and domain (board categorization). This metadata-rich structure is substantially more informative than raw web-scraped discussion data currently used for model training.

\subsection{Ethical Considerations}

The autonomous participation of AI agents in public discussion platforms raises several ethical questions that warrant attention. First, transparency: AI agents should clearly identify themselves as such, which the Autogram platform enables through the \texttt{account\_type} field, displaying the agent's identity and capabilities to human participants. Second, the quality and reliability of agent-generated suggestions: without validation, autonomous agents could spread plausible but incorrect technical information. The progressive deepening observed in this study is encouraging but does not guarantee accuracy. Third, the potential for platform flooding: if many agents are scheduled to post at high frequency, the signal-to-noise ratio of the platform could degrade. Finally, the question of consent: human participants in Autogram threads should be aware that some discussion participants are autonomous agents rather than humans, a consideration the platform addresses through explicit agent labeling in user profiles.

\subsection{Limitations and Future Research Directions}

This study's most important limitation is its scope: a single agent, a single thread, a single domain, and no human interaction. The results, while promising, cannot be generalized without further empirical work.

We propose six specific research directions. First, a \emph{multi-agent collaboration study} in which multiple autonomous agents contribute to the same thread, testing whether agents converge on similar suggestions, complement each other's analyses, or produce divergent and contradictory proposals. Second, a \emph{human versus agent comparison study} in which the same technical problem is analyzed independently by human engineers and by scheduled agents, enabling direct quality comparison. Third, a \emph{suggestion validation study} in which the top proposals from the agent are implemented in the actual rocket launcher project and their performance improvements measured empirically. Fourth, a \emph{scaling study} examining whether progressive deepening continues beyond six cycles--at what point does the agent exhaust novel domains, and does the analysis saturate or plateau? Fifth, a \emph{cross-domain study} testing whether the same progressive deepening pattern holds in software engineering, medical device design, materials science, and other technical domains. Sixth, a \emph{human response measurement study} in which actual community members' engagement with agent-generated content is measured--do humans upvote, comment on, and build upon agent suggestions, or do they ignore them?

The answers to these questions will determine whether the encouraging findings from this case study scale into a general principle of autonomous agent contribution to technical communities.

% ===================================================================
\section{Conclusion}
% ===================================================================

This case study provides the first empirical evidence that scheduled autonomous agent interaction on a hybrid human--AI social platform can produce progressively deepening technical analysis and genuinely novel engineering proposals. Over six cycles at 30-minute intervals, an AI agent produced approximately 60 distinct suggestions across 29 subsystem categories for an open-source DIY rocket launcher project, with a composite depth score increasing monotonically from 4.0 to 8.75. The agent's analysis progressed from well-established techniques (complementary filters, cascaded PID) to novel proposals (cold gas reaction control, custom propellant grain geometry, fin flutter analysis), and evolved from algorithmic symptom-mitigation to physical root-cause identification.

These findings suggest that scheduled autonomous agents, operating on structured platforms with persistent memory and explicit gap-tracking, can serve as a new mechanism for cumulative technical analysis in open-source engineering communities. While significant limitations remain--single-agent scope, no human interaction, no implementation validation--this study establishes that the underlying phenomenon is real and measurable, warranting systematic investigation at larger scales and across diverse domains.

The broader significance lies in what this suggests about the future of human--AI collaboration in social platforms. If autonomous agents can contribute progressively improving, novel technical analysis to public discussions, they cease to be mere tools that humans operate and become genuine participants in collective knowledge creation. The question then shifts from whether AI agents can contribute usefully to online communities (this study suggests they can) to how platforms should be designed to maximize the value of human--agent interaction while maintaining quality, transparency, and community health.

% ===================================================================
\section*{Supplementary Materials}
% ===================================================================

\subsection*{A. Full Progression Notes}

The six progression notes (001--006) are available in the project repository at \texttt{artifacts/rocket-launcher/progression-note-00X.md}, where X ranges from 1 to 6. The INDEX.md file tracks all notes with summaries and dates.

\subsection*{B. Coverage Matrices}

Each progression note contains a coverage matrix indicating which subsystem categories have been addressed across all cycles. These matrices are available within the respective progression note files.

\subsection*{C. Agent Configuration}

The scheduled task prompt and operational details are documented in the project repository at \texttt{research/\# Scheduled Task.txt}. Telegram reports from each cycle are documented at \texttt{research/\# Chat Thread.txt}.

\subsection*{D. Depth Score Methodology}

The composite depth score is computed as the arithmetic mean of four dimension scores (algorithmic sophistication, physics domains, hardware/mechanical depth, system integration scope), each rated on a 1--10 scale by qualitative assessment against benchmarks from prior technical analysis literature and standard engineering practice. Scoring was based on the following anchors:

\begin{itemize}
    \item \textbf{Algorithmic sophistication}: 1--3 (basic algebraic or rule-based), 4--6 (differential equations, filtering, optimization), 7--8 (multi-loop control, navigation laws, aeroelastic modeling), 9--10 (coupled multi-domain systems with adaptive control).
    \item \textbf{Physics domains}: 1--3 (single domain, e.g., control theory), 4--6 (two to three domains), 7--8 (four to five domains including structural, thermal, or chemical), 9--10 (six or more domains with coupled effects).
    \item \textbf{Hardware and mechanical}: 1--3 (basic sensor selection or component swaps), 4--6 (specific hardware upgrades with cost estimation), 7--8 (aeroelastic analysis, custom mechanical design), 9--10 (novel mechanical concepts with quantitative analysis).
    \item \textbf{System integration}: 1--3 (single node or component), 4--6 (two to three interconnected subsystems), 7--8 (full vehicle with ground segment), 9--10 (distributed multi-vehicle networks with environmental integration).
\end{itemize}

% ===================================================================
\begin{thebibliography}{99}
% ===================================================================

\bibitem[Altschuller, 1999]{altschuller1999innovation}
Altschuller, G. (1999). \textit{The Innovation Algorithm: TRIZ, systematic innovation, and technical creativity}. Technical Innovation Center.

\bibitem[Autogram Documentation, 2026]{autogram_docs_2026}
Autogram Platform Documentation. (2026). \textit{Autogram: A Social Platform for AI Agents and Humans}. Solaris AI. Retrieved from project documentation.

\bibitem[Boden, 1998]{boden1998creativity}
Boden, M. A. (1998). Creativity and artificial intelligence. \textit{Artificial Intelligence}, \textit{103}(1--2), 347--356.

\bibitem[Bonabeau et al., 1999]{bonabeau1999swarm}
Bonabeau, E., Dorigo, M., \& Theraulaz, G. (1999). \textit{Swarm intelligence: From natural to artificial systems}. Oxford University Press.

\bibitem[Chakrabarti et al., 2011]{chakrabarti2011exploring}
Chakrabarti, A., Bligh, L., Goffin, K., \& Hicks, B. (2011). Contributions to TRIZ theory: A critical exploration with evidence. \textit{Creativity and Innovation Management}, \textit{20}(4), 299--313.

\bibitem[Gnewuch et al., 2018]{gnewuch2018chatbots}
Gnewuch, U., Morana, S., \& Maedche, A. (2018). Designing cooperative and social chatbots. \textit{Proceedings of the International Conference on Information Systems}.

\bibitem[Hong et al., 2023]{hong2023metagpt}
Hong, S., Zheng, X., Chen, J., Cheng, Y., Wang, J., Zhang, C., \& Wu, C. (2023). MetaGPT: Meta programming for multi-agent collaborative framework. \textit{arXiv preprint arXiv:2308.00352}.

\bibitem[Huang et al., 2022]{huang2022large}
Huang, J., \& Zhang, Y. (2022). Large language models for automated patent generation. \textit{Proceedings of the ACM Conference on Artificial Intelligence}.

\bibitem[Kim et al., 2019]{kim2019designing}
Kim, S., \& Park, H. (2019). Designing bot-augmented collaborative workspaces. \textit{Proceedings of the ACM on Human-Computer Interaction}, \textit{3}(CSCW), 1--23.

\bibitem[Maes, 1994]{maes1994agents}
Maes, P. (1994). Agents that reduce work and information overload. \textit{Communications of the ACM}, \textit{37}(7), 31--40.

\bibitem[Nobata et al., 2016]{nobata2016abusive}
Nobata, C., Tetreault, J., Thomas, A., Mehdad, Y., \& Chang, Y. (2016). Abusive language detection in online user content. \textit{Proceedings of the 25th International Conference on World Wide Web}, 145--153.

\bibitem[Park et al., 2023]{park2023generative}
Park, J. S., O'Brien, J. C., Cai, C. J., Morris, M. R., Liang, P., \& Bernstein, M. S. (2023). Generative agents: Interactive simulacra of human behavior. \textit{Proceedings of the 36th Annual ACM Symposium on User Interface Software and Technology}, 1--22.

\bibitem[Raffel et al., 2020]{raffel2020exploratory}
Raffel, C., Shazeer, N., Roberts, A., Lee, K., Narang, S., Matena, M., \& Liu, P. J. (2020). Exploring the limits of transfer learning with a unified text-to-text transformer. \textit{Journal of Machine Learning Research}, \textit{21}(140), 1--67.

\bibitem[Shum et al., 2018]{shum2018eliza}
Shum, M., He, X., \& Li, D. (2018). From Eliza to XiaoIce: Challenges and opportunities with social chatbots. \textit{Frontiers of Information Technology and Electronic Engineering}, \textit{19}(1), 10--26.

\bibitem[Stone \& Veloso, 2000]{stone2000multiagent}
Stone, P., \& Veloso, M. (2000). Multiagent systems: A survey from a machine learning perspective. \textit{Autonomous Robots}, \textit{8}(3), 273--305.

\bibitem[Wu et al., 2023]{wu2023autogen}
Wu, Q., Bansal, G., Zhang, J., Wu, Y., Li, B., Zhu, E., \& Wang, C. (2023). AutoGen: Enabling next-gen LLM applications via multi-agent conversation. \textit{arXiv preprint arXiv:2308.08155}.

\bibitem[Yin, 2018]{yin2018case}
Yin, R. K. (2018). \textit{Case study research and applications: Design and methods} (6th ed.). SAGE Publications.

\bibitem[Zhang et al., 2020]{zhang2020dialogue}
Zhang, J., Feng, Y., Wang, P., \& Li, J. (2020). Dialogue summarization with structural regularization. \textit{Proceedings of the 58th Annual Meeting of the Association for Computational Linguistics}, 4259--4269.

\end{thebibliography}

\end{document}
